% !TEX root = ../../../main.tex
\section{Upper and lower sums}
\label{sec:analysis-i:upper-and-lower-sums}

% BEGIN TUFTE PLACEHOLDER: supplied sample, lines 958--985.
\begingroup\sloppy
% Replace this entire specimen block with reviewed course notes.
\label{template:sec:filehooks}

\index{file hooks|(}
If you create many documents using the \TL classes, it's easier to
store your customizations in a separate file instead of copying them into the
preamble of each document.  The \TL classes provide three file hooks:
\docfilehook{tufte-common-local.tex}{common}, \docfilehook{tufte-book-local.tex}{book}, and
\docfilehook{tufte-handout-local.tex}{handout}.

\begin{description}
  \item[\docfilehook{tufte-common-local.tex}{common}]
    If this file exists, it will be loaded by all of the \TL document
    classes just prior to any document-class-specific code.  If your
    customizations or code should be included in both the book and handout
    classes, use this file hook.
  \item[\docfilehook{tufte-book-local.tex}{book}] 
    If this file exists, it will be loaded after all of the common and
    book-specific code has been read.  If your customizations apply only to the
    book class, use this file hook.
  \item[\docfilehook{tufte-common-handout.tex}{handout}] 
    If this file exists, it will be loaded after all of the common and
    handout-specific code has been read.  If your customizations apply only to
    the handout class, use this file hook.
\end{description}

\index{file hooks|)}

\lipsum[1]
\par\endgroup
% END TUFTE PLACEHOLDER
